\documentclass[../../main.tex]{subfiles}
\begin{document}

{\bf [解]}

数学的帰納法により題意の不等式
\begin{align}
  \sum_{j=1}^{n}(x_j-y_j)^2 \le \sum_{j=1}^{n}(x_j-z_j)^2 \label{eq:1}
\end{align}
を示す．

まず$n=2$の場合を考える．
数列$\{x_1,x_2\}$および$\{y_1,y_2\}$と，$\{y_n\}$を並び替えた$\{z_1,z_2\}$について考える．並び替えのパターンとしては$(y_1,y_2)=(z_1,z_2),(z_2,z_1)$のふた通りである．前者の時は\cref{eq:1}で等号が成立するので良い．そこで後者の場合を考える．
右辺と左辺の差を計算すると
\begin{align}
  (\text{右辺})-(\text{左辺})
  &= \sum_{k=1}^2 (x_k - z_k)^2 - \sum_{k=1}^2 (x_k - y_k)^2 \\
  &= (-2x_1 y_2 - 2x_2 y_1) - (-2x_1 y_1 - 2x_2 y_2) \\
  &= 2(x_1 - x_2)(y_1 - y_2) \ge 0 \quad (\because x_1 \ge x_2, \, y_1 \ge y_2)
\end{align}
より\cref{eq:1}は成立する．
以上より全ての場合で\cref{eq:1}が成立する．

そこで次に$n=k\ge 2$での成立を仮定して$n=k+1$の場合を考える．
数列$\{x_n\},\{y_n\}$が
\begin{align}
  \begin{cases}
    x_1 \ge x_2 \ge x_3 \ge \cdots \ge x_{k} \ge x_{k+1} \\
    y_1 \ge y_2 \ge y_3 \ge \cdots \ge y_{k} \ge y_{k+1}
  \end{cases}
\end{align}
を満たすとする．$z_{k+1}$の値に応じて場合分けする．

\subparagraph{$1^{\circ}$ $z_{k+1} = y_{k+1}$ の時}

この時は$(x_1,x_2,\cdots,x_{k}), (y_1,y_2,\cdots,y_{k})$に仮定を適用できるから
$n=k+1$の\cref{eq:1}の右辺と左辺の差は
\begin{align}
  (\text{右辺})-(\text{左辺})
  &\sum_{i=1}^{k+1} (x_i - z_i)^2 - \sum_{i=1}^{k+1} (x_i - y_i)^2 \\
  &= \sum_{i=1}^k (x_i - z_i)^2 + (x_{k+1}-y_{k+1})^2 - \sum_{i=1}^k (x_i - y_i)^2 -  (x_{k+1}-y_{k+1})^2 \\
  &= \sum_{i=1}^k (x_i - z_i)^2 - \sum_{i=1}^k (x_i - y_i)^2 \ge 0 \quad (\because \text{仮定})
\end{align}
だから $n = k+1$ でも\cref{eq:1}は成立する．

\subparagraph{$2^{\circ}$ $z_{k+1} \neq y_{k+1}$ の時}

$y_{k+1}$ は $z_{k+1}$ 以外のある $z_l$ ($l \le k$) に等しく，
$z_{k+1}$ は $y_{k+1}$ 以外のある $y_m$ ($m \le k$) に等しい．
ここで，次のような数列 $\{z'_n\}$ をかんがえる．

\begin{itemize}
  \item $z'_{k+1} = z_l = y_{k+1}$ \\
  \item $z'_l = z_{k+1} = y_m$ \\
  \item $j \neq l, k+1$ の時 $z'_j = z_j$
\end{itemize}

% \begin{table}[htb]
%   \centering
%   \begin{tabular}{l|l}
%     $z'_{k+1} = z_l = y_{k+1}$ & $z_1, z_2, \dots, z_l, \dots, z_k, z_{k+1}$ \\
%     $z'_l = z_{k+1} = y_m$ & $z_1, z_2, \dots, z_{k+1}, \dots, z_k, z_l$ \\
%     $j \neq l, k+1$ の時 $z'_j = z_j$ & $y_{k+1}, \dots, y_m$
%   \end{tabular}
%   \caption{数列 $\{z'_n\}$ の定義}
% \end{table}

この時$\{y_n\}$ と $\{z'_n\}$ において，
$y_{k+1} = z'_{k+1}$ であり，
$1 \le n \le k$ の時 $\{z'_n\}$ は $\{y_n\}$ をならびかえたものだから
仮定から
\begin{align}
  \sum_{i=1}^{k} (x_i - z'_i)^2 + (x_{k+1} - z'_{k+1})^2 \ge \sum_{i=1}^{k} (x_i - y_i)^2 + (x_{k+1} - y_{k+1})^2 \\
  \sum_{i=1}^{k+1} (x_i - z'_i)^2 \ge \sum_{i=1}^{k+1} (x_i - y_i)^2 \label{eq:2}
\end{align}
が成り立つ．

一方で$\{z_n\}$ と $\{z'_n\}$ を比較すると
\begin{align}
  &\sum_{i=1}^{k+1} \Bigl[ (x_i - z_i)^2 - (x_i - z'_i)^2 \Bigr] \\
  &= \Bigl\{ (x_l - z_l)^2 - (x_l - z'_l)^2 \Bigr\} + \Bigl\{ (x_{k+1} - z_{k+1})^2 - (x_{k+1} - z'_{k+1})^2 \Bigr\} \\
  &= \Bigl\{ (x_l - y_{k+1})^2 - (x_l - y_m)^2 \Bigr\} + \Bigl\{ (x_{k+1} - y_m)^2 - (x_{k+1} - y_{k+1})^2 \Bigr\} \\
  &= -2x_l y_{k+1} + 2x_l y_m - 2x_{k+1} y_m + 2x_{k+1} y_{k+1} \\
  &= 2(x_{k+1} - x_l)(y_{k+1} - y_m) \ge 0 \quad (\because y_m \le y_{k+1}, \, x_l \le x_{k+1}) \\
  \therefore
  \sum_{i=1}^{k+1} (x_i - z_i)^2 \ge \sum_{i=1}^{k+1}(x_i - z'_i)^2 \label{eq:3}
\end{align}
である．

\cref{eq:2,eq:3}より
\begin{align}
  \sum_{i=1}^{k+1} (x_i - z_i)^2 \ge \sum_{i=1}^{k+1}(x_i - z'_i)^2 \ge \sum_{i=1}^{k+1} (x_i - y_i)^2 \\
  \sum_{i=1}^{k+1} (x_i - z_i)^2 \ge \sum_{i=1}^{k+1} (x_i - y_i)^2
\end{align}
となり，$n = k+1$ でも\cref{eq:1}は成立する．
\casesend

以上より数学的帰納法により任意の$n$について\cref{eq:1}が成立する．
よって題意は示された．

{\bf [別解]}

数列 $\{z_n\}$ に対し，まず $z_1$ と $z_l$ ($l = 2, 3, \dots, n$) を順に比較し，$z_1 < z_l$ の時に限り $z_1$ と $z_l$ を交換し，交換後の $z_1$ に対し，交換後の $z_l$ の次の項以下について同様の操作を続ける．次に交換後の $z_2$ と $z_l$ ($3 \le l \le n$) について上と同様の作業をくり返す．この操作が全て終了すると，$\{z_n\}$ は $\{y_n\}$ となる．

次に各操作において $z_p < z_q$ ($p < q$) の時 $z_p$ と $z_q$ を交換後 $z'_p = z_q, z'_q = z_p$ となったとすると，
\begin{align}
  &(x_p - z'_p)^2 + (x_q - z'_q)^2 - \Bigl\{ (x_p - z_p)^2 + (x_q - z_q)^2 \Bigr\} \\
  &= 2(x_p - x_q)(z_p - z_q) \le 0
\end{align}
だから，
\begin{align}
  (x_p - z'_p)^2 + (x_q - z'_q)^2 \le (x_p - z_p)^2 + (x_q - z_q)^2
\end{align}
したがって，$P_n = \sum_{k=1}^n (x_k - z_k)^2$ の操作後の値を $P'_n$ とすれば $P'_n \le P_n$ となる．全ての操作がおわった時のことをかんがえて
\begin{align}
  \sum_{k=1}^n (x_k - y_k)^2 \le \sum_{k=1}^n (x_k - z_k)^2
\end{align}
である．

\end{document}
